\pagestyle{plain}

            %%% ШРИФТ, АБЗАЦ, ПОЛЯ %%%

\usepackage{mathtext}
\usepackage[utf8]{inputenc}
\usepackage{cmap}                   % для кодировки шрифтов в pdf
\usepackage[pdftex]{lscape}         % Для ориентации страниц
\usepackage[14pt]{extsizes}         % Размер шрифта 14pt.


\setlength{\parindent}{2.5em}       % Размер отступа красной строки в 5 символов
\righthyphenmin=2
\usepackage{ulem}                   % Подчеркивания
\usepackage[usenames, dvipsnames]{color} % цвета

\usepackage[table,xcdraw]{xcolor}


%\newcommand{\basegostsectionfont}{\fontsize{14pt}{16pt}\selectfont\bfseries}


            %%% РУБРИКАЦИЯ %%%
            
\usepackage{misccorr}       % В заголовках появляется точка, но при ссылке на них ее нет 

\clubpenalty=10000      % Висячие строки.
\widowpenalty=10000
\tolerance=2000         % Терпимость к жидким строкам.
\righthyphenmin=2       % Минимальное число символов при переносе.

\usepackage{fancyhdr} % пакет для установки колонтитулов
\pagestyle{fancy} % смена стиля оформления страниц
\fancyhf{} % очистка текущих значений
\fancyhead[C]{\normalsize\thepage} % установка верхнего колонтитула
\renewcommand{\headrulewidth}{0pt} % убрать разделительную линию

\usepackage{titlesec}

% Оформление разделов
\makeatletter
    %\renewcommand{\l@section}{\@dottedtocline{1}{0.8cm}{0.8cm}}
    \renewcommand{\thesection}{\arabic{section}}
    \renewcommand{\section}{\@startsection{section}{1}{1.25cm}{-3.5ex plus -1ex minus -.2ex}{2.3ex plus.2ex}{\raggedright\normalfont\large\bfseries}}
\makeatother

% Оформление подразделов
\makeatletter
    %\renewcommand{\l@subsection}{\@dottedtocline{2}{0.8cm}{0.8cm}}
    \renewcommand{\thesubsection}{\arabic{section}.\arabic{subsection}}
    \renewcommand{\subsection}{\@startsection{subsection}{2}{1.25cm}{-3.5ex plus -1ex minus -.2ex}{2.3ex plus.2ex}{\raggedright\normalfont\large\bfseries}}
\makeatother

% Оформление под-подразделов
\makeatletter
    \renewcommand{\thesubsubsection}{\arabic{section}.\arabic{subsection}.\arabic{subsubsection}}
    \renewcommand{\subsubsection}{\@startsection{subsubsection}{3}{1.25cm}{-3.5ex plus -1ex minus -.2ex}{2.3ex plus.2ex}{\raggedright\normalfont\bfseries}}
\makeatother

% Оформление параграфов
\makeatletter
    \renewcommand{\theparagraph}{\arabic{section}.\arabic{subsection}.\arabic{subsubsection}.\arabic{paragraph}}
    \renewcommand{\paragraph}{\@startsection{paragraph}{4}{1.25cm}{-3.5ex plus -1ex minus -.2ex}{2.3ex plus.2ex}{\raggedright\normalfont\bfseries}}
\makeatother

\makeatletter
\renewcommand{\l@section}{\@dottedtocline{1}{0em}{1.25em}}
\renewcommand{\l@subsection}{\@dottedtocline{2}{1.25em}{1.75em}}
\renewcommand{\l@subsubsection}{\@dottedtocline{3}{2.75em}{2.6em}}
\makeatother

            %%% ФОРМУЛЫ %%%
% Запятая в десятичных дробях.            
\usepackage{icomma}
% AMS-пакеты.
\usepackage{amssymb, amsfonts, amsmath, amsthm, amstext, bm}
% Нормальные интегралы.
\usepackage{esint}
\renewcommand{\theequation}{\arabic{subsection}.\arabic{equation}}
\numberwithin{equation}{section}    % Нумерация формул по секциям.

\newcommand{\integral}[4]
{
\displaystyle \int\limits_{#1}^{#2}#3\,\mathrm{d}#4
}

\newcommand{\diff}[2]
{
\cfrac{\mathrm{d}#1}{\mathrm{d}#2}
}


            %%% РИСУНКИ, ТАБЛИЦЫ, СПИСКИ %%%
            
\usepackage{graphicx}               % Пакет для работы с графикой
\usepackage{longtable}              % Длинные таблицы
\usepackage{multirow}               % Для оформления сложных таблиц
\usepackage{array}
\usepackage{float}
\renewcommand{\labelenumi}{\theenumi.}
\newenvironment{tscript}            % Окружение для ввода обозначений (два столбца).
{\begin{tabular}{p{5mm}p{\dimexpr 1\linewidth-7mm-2\tabcolsep}}}{\end{tabular}}

% объявляем новую команду для переноса строки внутри ячейки таблицы
\newcommand{\specialcell}[2][c]{%
  \begin{tabular}[#1]{@{}c@{}}#2\end{tabular}}

            %%% ПОДПИСИ, ССЫЛКИ %%%
            
\usepackage{cleveref}
%\newcommand{\crefrangeconjunction}{--}
\crefrangeformat{equation}{#3(#1)#4--#5(#2)#6}
            
\usepackage{caption}
\captionsetup{%
singlelinecheck=off,                % Многострочные подписи, например у таблиц
skip=2pt,                           % Вертикальная отбивка между подписью и содержимым рисунка или таблицы
tableposition=top,                  % Подпись таблиц сверху
labelsep=period,                    % Разделитель - точка и пробел.
format=plain, 
font=small, 
labelfont=small,
}
\captionsetup[figure]{justification=centering}
\captionsetup[table]{justification=raggedright}

\usepackage{chngcntr}
\counterwithin{figure}{section}
\counterwithin{table}{section}
\usepackage{comment}

% Греческие символы прямым шрифтом.
%\def\alpha{{\text{\usefont{U}{psy}{m}{n}\selectfont\symbol{97}}}}
%\def\beta{{\text{\usefont{U}{psy}{m}{n}\selectfont\symbol{98}}}}
%\def\chi{{\text{\usefont{U}{psy}{m}{n}\selectfont\symbol{99}}}}
%\def\delta{{\text{\usefont{U}{psy}{m}{n}\selectfont\symbol{100}}}}
%\def\varepsilon{{\text{\usefont{U}{psy}{m}{n}\selectfont\symbol{101}}}}
%\def\phi{{\text{\usefont{U}{psy}{m}{n}\selectfont\symbol{102}}}}
%\def\gamma{{\text{\usefont{U}{psy}{m}{n}\selectfont\symbol{103}}}}
%\def\eta{{\text{\usefont{U}{psy}{m}{n}\selectfont\symbol{104}}}}
%\def\iota{{\text{\usefont{U}{psy}{m}{n}\selectfont\symbol{105}}}}
%\def\varphi{{\text{\usefont{U}{psy}{m}{n}\selectfont\symbol{106}}}}
%\def\kappa{{\text{\usefont{U}{psy}{m}{n}\selectfont\symbol{107}}}}
%\def\lambda{{\text{\usefont{U}{psy}{m}{n}\selectfont\symbol{108}}}}
%\def\mu{{\text{\usefont{U}{psy}{m}{n}\selectfont\symbol{109}}}}
%\def\nu{{\text{\usefont{U}{psy}{m}{n}\selectfont\symbol{110}}}}
%\def\omicron{{\text{\usefont{U}{psy}{m}{n}\selectfont\symbol{111}}}}
%\def\pi{{\text{\usefont{U}{psy}{m}{n}\selectfont\symbol{112}}}}
%\def\theta{{\text{\usefont{U}{psy}{m}{n}\selectfont\symbol{113}}}}
%\def\rho{{\text{\usefont{U}{psy}{m}{n}\selectfont\symbol{114}}}}
%\def\sigma{{\text{\usefont{U}{psy}{m}{n}\selectfont\symbol{115}}}}
%\def\tau{{\text{\usefont{U}{psy}{m}{n}\selectfont\symbol{116}}}}
%\def\upsilon{{\text{\usefont{U}{psy}{m}{n}\selectfont\symbol{117}}}}
%\def\varpi{{\text{\usefont{U}{psy}{m}{n}\selectfont\symbol{118}}}}
%\def\omega{{\text{\usefont{U}{psy}{m}{n}\selectfont\symbol{119}}}}
%\def\xi{{\text{\usefont{U}{psy}{m}{n}\selectfont\symbol{120}}}}
%\def\psi{{\text{\usefont{U}{psy}{m}{n}\selectfont\symbol{121}}}}
%\def\zeta{{\text{\usefont{U}{psy}{m}{n}\selectfont\symbol{122}}}}
%\def\epsilon{\varepsilon}
%\DeclareMathSymbol{\int}{\mathop}{Symbol}{242}
%\def\intup{\mbox{\usefont{U}{psy}{m}{n}\selectfont\symbol{243}}}
%\def\intdn{\mbox{\usefont{U}{psy}{m}{n}\selectfont\symbol{245}}}
%\def\int{\mathop{\intup\atop\intdn}}
%\def\int{\mathop{\raisebox{-.7ex}{\usefont{U}{psy}{m}{n}\fontsize{16pt}{16pt}\selectfont\symbol{242}}}}
%\def\oint{\mathop{\lefteqn{\circ}\raisebox{-1ex}{\usefont{U}{psy}{m}{n}\fontsize{16pt}{16pt}\selectfont\symbol{242}}}}